\documentclass[aspectratio=169]{beamer}

% Clean, modern theme setup
\usetheme{Boadilla}
\usecolortheme{seahorse} % Soft, clean color palette
\usefonttheme{professionalfonts} % Cleaner math rendering

\usepackage{amsmath}

\title{Pion Decay Rates in a 100 cm Tracking Volume}
\subtitle{Kinetic Energy Impact on Tracking Efficiency}
\author{}
\date{}

\begin{document}

\begin{frame}
    \titlepage
\end{frame}

\begin{frame}{1. The Decay Formula}
    In the lab frame, the mean decay length $\lambda$ is dilated by relativity:
    \[ \lambda = \beta\gamma c\tau \]
    Where:
    \begin{itemize}
        \item $c\tau \approx 7.80 \text{ m}$ (the proper decay length of a pion, using its rest lifetime $\tau \approx 26.03 \text{ ns}$)
        \item The probability of the pion decaying before reaching distance $L$ is:
        \[ P_{\text{decay}} = 1 - e^{-L/\lambda} = 1 - e^{-L/(\beta\gamma c\tau)} \]
    \end{itemize}
\end{frame}

\begin{frame}{2. Decay Chances Over 100 cm ($L = 1.0 \text{ m}$)}
    Using the two kinetic energy values:
    \vspace{0.5em}
    
    \textbf{1 GeV Kinetic Energy ($\beta\gamma \approx 8.16$)}
    \begin{itemize}
        \item Mean decay length: 
        \[ \lambda = 8.16 \times 7.80 \text{ m} \approx 63.6 \text{ m} \]
        \item Decay probability over 100 cm: 
        \[ P_{\text{decay}} = 1 - e^{-1.0/63.6} = 1 - e^{-0.0157} \approx 1.56\% \]
        \item \textit{(\textasciitilde 98.4\% survival chance)}
    \end{itemize}
\end{frame}

\begin{frame}{2. Decay Chances Over 100 cm (Continued)}
    \textbf{0.1 GeV Kinetic Energy ($\beta\gamma \approx 1.40$)}
    \begin{itemize}
        \item Mean decay length: 
        \[ \lambda = 1.40 \times 7.80 \text{ m} \approx 10.9 \text{ m} \]
        \item Decay probability over 100 cm: 
        \[ P_{\text{decay}} = 1 - e^{-1.0/10.9} = 1 - e^{-0.0917} \approx 8.76\% \]
        \item \textit{(\textasciitilde 91.2\% survival chance)}
    \end{itemize}
\end{frame}

\begin{frame}{3. Decay Chances Over 1 cm ($L = 0.01 \text{ m}$)}
    \textbf{1 GeV Kinetic Energy ($\beta\gamma \approx 8.16$)}
    \begin{itemize}
        \item Decay probability over 1 cm:
        \[ P_{\text{decay}} = 1 - e^{-0.01/63.6} = 1 - e^{-1.57\times10^{-4}} \approx 0.0157\% \]
        \item \textit{(\textasciitilde 99.984\% survival chance)}
    \end{itemize}
    \vspace{0.6em}
    \textbf{0.1 GeV Kinetic Energy ($\beta\gamma \approx 1.40$)}
    \begin{itemize}
        \item Decay probability over 1 cm:
        \[ P_{\text{decay}} = 1 - e^{-0.01/10.9} = 1 - e^{-9.17\times10^{-4}} \approx 0.0917\% \]
        \item \textit{(\textasciitilde 99.908\% survival chance)}
    \end{itemize}
\end{frame}

\begin{frame}{Impact on Tracking Efficiency}
    In a 100 cm tracking volume, approximately \textbf{8.8\%} of 0.1 GeV pions will decay into a muon and a neutrino ($\pi^+ \to \mu^+\nu_\mu$) before traversing the entire detector, causing a kinked or broken track, compared to only \textbf{1.6\%} of 1 GeV pions.
\end{frame}

\end{document}